\section{Rees families and Laurent bounds}\label{sec:rees}

We first record a reduction that can be applied to the original pair,
before augmentation normalization or any choice of filtration.

\begin{lemma}\label{lem:coefficient-field}
The given isomorphism $\Phi_{\mathrm{raw}}:U_k(L)\simeq U_k(M)$
is the scalar extension of an isomorphism
$\Phi_{K,\mathrm{raw}}:U_{K}(L_K)\simeq U_{K}(M_K)$ over a
subfield $K\subseteq k$ finitely generated over $\QQ$. If $L$
is solvable, so is $L_K$. The field $K$ embeds in $\CC$.
\end{lemma}

\begin{proof}
Choose finite ordered bases $x_1,\ldots,x_a$ and $y_1,\ldots,y_b$.
Take all coefficients in their two bracket tables and in the finite PBW
expansions of $\Phi_{\mathrm{raw}}(x_i)$ and $\Phi_{\mathrm{raw}}^{-1}(y_j)$. If the corresponding
supports are $S_i$ and $T_j$, this uses at most
\[
 a^3+b^3+\sum_i|S_i|+\sum_j|T_j|
\]
scalars. Let $K$ be the field they generate over $\QQ$. The same
bracket tables define Lie algebras over $K$: their identities are
reflected by the inclusion into $k$. Their scalar extensions recover
the original algebras.

The evaluation maps from the smaller enveloping algebras to the original
ones are injective by PBW, since they preserve the indexed ordered
monomials. The displayed generator images therefore satisfy the defining
relations over $K$, and define algebra maps in both directions.
Their inverse identities are likewise reflected by evaluation. This
descends the specified $\Phi_{\mathrm{raw}}$, including its inverse. Solvability
descends through the injective Lie map from $L_K$ to
$L$ with scalars restricted.
There is no need to descend an infinite list of enveloping elements:
the finitely many generator images determine each algebra map.
Including the inverse images ensures that the resulting map over $K$
is an equivalence, rather than merely a homomorphism whose scalar
extension is the desired forward map.

Finally, embed a finite transcendence basis of $K/\QQ$ as an
algebraically independent complex family. Extend that embedding across
the remaining finite algebraic extension into $\CC$.
\end{proof}

We may apply the preceding intrinsic argument over $K$. A Lie
isomorphism obtained there will extend to $k$; no embedding of $k$
in $\CC$ is needed. We henceforth rename the $K$-models $L,M$ and their equivalence
$\Phi_{\mathrm{raw}}$. Normalize this descended map as in \cref{lem:augmentation}
and again denote it by $\Phi$. The character used in that
normalization belongs to $K$: its values are constant PBW
coefficients already included in $K$. All adapted choices
and the marking $\varphi_0$ are now made over $K$.

\subsection{The polynomial families}

Choose bases $x_1,\ldots,x_n$ and $y_1,\ldots,y_n$ adapted to the
recovered filtrations and matched by the recovered graded isomorphism $\varphi_0$.
Write $w_i\geq0$ for the common weights. With
$\mathcal F^qU(M)$ denoting the span of PBW monomials of
weight at least $q$, the recovered marking gives
\begin{equation}\label{eq:rees-leading}
 \Phi(x_i)-y_i\in\mathcal F^{w_i+1}U(M),\qquad
 \Phi^{-1}(y_i)-x_i\in\mathcal F^{w_i+1}U(L).
\end{equation}
These follow from \cref{prop:graded-mark}, rather than being additional assumptions on the
original isomorphism. The bracket coefficients satisfy
$c_{ij}^p\ne0\Rightarrow w_p\geq w_i+w_j$, and similarly for
$M$. Distinguish the weight
$w(\alpha)=\sum_iw_i\alpha_i$ from the ordinary degree
$|\alpha|=\sum_i\alpha_i$.
Concretely, choose vector-space sections of the finitely many graded
pieces and lift a homogeneous basis through them. Use $\varphi_0$
to choose the corresponding target basis. These are linear
splittings, with no assertion that either splitting preserves brackets.
The filtration is the span of basis vectors of weight at least
its index. Thus $[F^pL,F^qL]\subseteq F^{p+q}L$ gives the
displayed support condition. The degree-zero part and the
zero-dimensional case are included.

Define the free $K[t]$-module $L_t$ by
\begin{equation}\label{eq:rees-bracket}
 [x_i,x_j]_t=\sum_{p:c_{ij}^p\ne0}
 c_{ij}^p t^{w_p-w_i-w_j}x_p.
\end{equation}
All exponents are nonnegative. After extending to $K((t))$, the
diagonal map $\delta_L(x_i)=t^{-w_i}x_i$ preserves this bracket:
\[
 \delta_L([x_i,x_j]_t)
 =\sum_pc_{ij}^pt^{-w_i-w_j}x_p
 =[\delta_L(x_i),\delta_L(x_j)].
\]
Its restriction to the polynomial module is injective. Hence Jacobi
is reflected from the original scalar-extended Lie algebra.
Construct $M_t$ similarly. At $t=0$ exactly the
equal-weight terms survive, giving the common graded Lie algebra
$g_0$; at $t=1$ the original algebras return.

We use PBW over $K[t]$, and later over $K[t,h]$. For these free
Lie modules, ordered monomials span by commutator rewriting. Mapping
to the enveloping algebra over the fraction field, field PBW shows
that every finite relation between these monomials has zero
coefficients. Thus they are a basis over the polynomial ring, and
localization is injective. This argument does not presuppose the
injectivity it establishes. Symmetrization is triangular with
diagonal one for ordinary degree, so it and its inverse preserve
the ordinary degree filtration.

\begin{proposition}\label{prop:rees-isomorphism}
There is a $K[t]$-algebra isomorphism
\[
 \Phi_t:U_{K[t]}(L_t)\simeq U_{K[t]}(M_t)
\]
whose zero specialization is the identity under the fixed marking.
If $\Phi(x_i)=\sum_\alpha a_{i\alpha}y^\alpha$, then
\begin{equation}\label{eq:rees-map}
 \Phi_t(x_i)=\sum_\alpha
 a_{i\alpha}t^{w(\alpha)-w_i}y^\alpha.
\end{equation}
\end{proposition}

\begin{proof}
Filteredness makes all nonzero terms in \eqref{eq:rees-map}
polynomial. After Laurent scalar extension the formula is
$U(\delta_M)^{-1}\Phi U(\delta_L)$, so it preserves the relations.
Injective localization reflects these relations to the polynomial
envelope. Applying the same construction to $\Phi^{-1}$ gives a
polynomial inverse; the inverse identities can again be checked
after localization. By \eqref{eq:rees-leading},
$\Phi_t(x_i)=y_i+tQ_i(t)$ for an actual polynomial element
$Q_i(t)\in U_{K[t]}(M_t)$. The zero specialization fixes
every marked generator and therefore the whole envelope.
\end{proof}

Let $W$ be the common coordinate vector space and $A_K=S_K(W)$.
In symmetric PBW coordinates put
\[
 \Psi_t=\operatorname{sym}_{M_t}^{-1}\Phi_t\operatorname{sym}_{L_t},
 \qquad \Psi_t=\sum_{r\geq0}t^r\Psi_r,\qquad \Psi_0=\Id.
\]
Here coefficient extraction defines $\Psi_r\in\End_K(A_K)$;
each fixed polynomial input has a polynomial output in $t$.
The map intertwines the two transported enveloping products,
not necessarily ordinary commutative multiplication.

\subsection{Homogenization and the degree estimate}

Introduce an independent parameter $h$, with
$U_h(L_t)$ defined by
$xy-yx=h[x,y]_t$. This is the envelope of the bracket
multiplied by $h$ and has the same PBW basis. Inverting
$h$ gives the algebra isomorphism
\[
 \theta_{h,L}:U_h(L_t)[h^{-1}]
 \longrightarrow U(L_t)[h,h^{-1}],
 \qquad x_i\longmapsto h x_i.
\]
Both sides of its defining relation map to $h^2[x_i,x_j]_t$.
Consequently
\begin{equation}\label{eq:rees-homogenized}
 \Phi_{t,h}=\theta_{h,M}^{-1}\Phi_t\theta_{h,L},
 \qquad
 \Phi_{t,h}(x_i)=\sum_\alpha
 a_{i\alpha}t^{w(\alpha)-w_i}h^{1-|\alpha|}y^\alpha.
\end{equation}
This construction initially concerns finite polynomials. Its scalar
extension to $K((h))[t]$ preserves the outer variable $t$.

Symmetrization of $n$ generators scales by exactly $h^n$.
Thus the block of $\Psi_r$ from $S^n(W)$ to $S^m(W)$ is
multiplied under \eqref{eq:rees-homogenized} by
\begin{equation}\label{eq:rees-block}
 h^{n-m}.
\end{equation}
Indeed, every ordered product in the normalized permutation average
has $n$ factors, and hence acquires the same input multiplier.
Inverse scaling contributes $h^{-m}$ to its output component.
This establishes the exact identity in symmetric coordinates;
an estimate in ordered PBW coordinates alone would not identify
the operator used below.
Choose $D\geq1$ bounding the ordinary PBW degrees of all terms in
the images of basis generators under both $\Phi$ and $\Phi^{-1}$.
The finite bases and finite PBW supports ensure such a $D$ exists.

\begin{lemma}\label{lem:uniform-excess}
For every $n,r\geq0$,
\[
 \Psi_r(S^n(W))\subseteq\bigoplus_{m\leq n+r(D-1)}S^m(W).
\]
The same bound holds for the coefficients of the inverse comparison.
\end{lemma}

\begin{proof}
Work in the actual algebra $B_t=U_{K[t]}(M_t)$,
with its ordinary PBW filtration $B_{t,\leq d}$.
This filtration is multiplicative over $K[t]$.
The corrections in $\Phi_t(x_i)=y_i+tQ_i(t)$ belong to
$B_{t,\leq D}$: subtracting a linear generator and reweighting
its coefficients by powers of $t$ cannot increase this degree bound.

Introduce an auxiliary central variable $z$ and consider
$y_i+zQ_i(t)$ in $B_t[z]$. In a product of $n$ factors, a
coefficient of $z^s$ selects $s$ corrections and has degree at most
\[
 (n-s)+sD=n+s(D-1).
\]
Apply this estimate to the finite symmetrization average of the
input monomial. Inverse target symmetrization is $K[t]$-linear
and degree-nonincreasing, so it preserves this bound.

Finally evaluate $z=t$. The contribution of a $z^s$ coefficient
to $t^r$ has residual polynomial degree $r-s$, and therefore
$s\leq r$. Its spatial degree is at most $n+r(D-1)$.
This proves the assertion; interchange source and target for
the inverse. Throughout, multiplication remains the genuine
$t$-dependent multiplication in $B_t$.
\end{proof}

The estimate bounds degree \emph{excess} uniformly over every input
degree at fixed $r$. Its dependence on $r$ is essential.

\subsection{The completed operator}

Extend coefficients along a fixed embedding $K\hookrightarrow\CC$,
and now write $A=S_\CC(W_\CC)$ and $C=D-1$.
For any complex vector space $X$ set
\[
 X((h))=
 \left\{\sum_{d\geq B}h^dx_d:
 B\in\mathbb Z,\ x_d\in X\right\},\qquad
 E=\CC((h)),\quad R=E[[t]].
\]
There is no finite-dimensionality restriction on the coefficient
span. For example $\sum_{j\geq0}h^jx^j$ belongs to
$\CC[x]((h))$ but not to $E[x]$.
Our function space is
$\mathcal A=A((h))[[t]]$: each $t$-coefficient has its own
Laurent lower bound.

\begin{proposition}\label{prop:Laurent-operator-bound}
The finite comparison \eqref{eq:rees-homogenized} defines an element
\[
 T=\Id+\sum_{r\geq1}t^rT_r
 \in\mathcal E:=\End_\CC(A)((h))[[t]],
 \qquad \operatorname{ord}_h T_r\geq-rC,
\]
where the bound holds on the whole coefficient endomorphism.
It acts invertibly and $R$-linearly on $\mathcal A$, intertwining
the two completed homogenized PBW products.
\end{proposition}

\begin{proof}
Define $T_{r,d}$ on $S^n(W_\CC)$ to be the output component
of $\Psi_r$ in degree $n-d$, interpreted as zero for $n-d<0$.
This is an endomorphism of $A$, since each input has finitely
many homogeneous components. The uniform-excess lemma gives
$T_{r,d}=0$ whenever $d<-rC$. Hence
$T_r=\sum_{d\geq-rC}h^dT_{r,d}$ is a Laurent series of
endomorphisms. Formula \eqref{eq:rees-block} identifies it with
the actual homogenized comparison. On an input of degree at most
$n$, its Laurent support is contained in $[-rC,n]$.

For $F=\sum_{d\geq B}h^dF_d$ and
$x=\sum_{j\geq b}h^jx_j$, define
\[
 [h^q](Fx)=\sum_{d=B}^{q-b}F_d(x_{q-d}).
\]
This is finite and has output lower bound $B+b$.
Outer convolution gives
$[t^r](Tx)=\sum_{p=0}^rT_p(x_{r-p})$.
If $x_s$ has lower bound $b_s$, the latter has lower bound
$\min_{0\leq p\leq r}(b_{r-p}-pC)$.
These formulas give the asserted scalar-linear action and preserve
operator composition. With $G=T-\Id$, the inverse
$\sum_{s\geq0}(-G)^s$ exists because $G$ has positive $t$-order.

The homogenized PBW products have nonnegative $h$-exponents
as whole binary cochains. More explicitly, a multiplication block
from $S^a(W)\otimes S^b(W)$ to $S^m(W)$ acquires the factor
$h^{a+b-m}$ under homogenization, and the ordinary PBW filtration
gives $m\leq a+b$. Thus each coefficient is a binary cochain on
all of $A$, with lower $h$-bound zero, and the products extend
by the same convolution. On finite polynomial inputs,
intertwining follows from the algebra isomorphism
\eqref{eq:rees-homogenized}. Extracting coefficients gives equality
of the binary coefficient maps; convolution then proves it for
all inputs in $\mathcal A$.
\end{proof}

This is an explicit action of coefficient series, not an
identification with all linear or continuous maps on $\mathcal A$.
The action is faithful: evaluate it on constant inputs $a\in A$,
then extract first the $t$-coefficient and then the
$h$-coefficient to recover each endomorphism applied to $a$.
The same argument detects equality of multilinear coefficient
families. It does not require constant inputs to span the completed
module algebraically over $R$.
Likewise spatial scaling is only used on finite polynomials:
substituting $x\mapsto h^{-1}x$ in
$\sum_{j\geq0}h^jx^{2j}$ would produce an inadmissible
series with no Laurent lower bound. The uniform estimate is what
constructs the completed comparison. No lower bound uniform in
$t$-degree is needed; $\sum_{r\geq0}t^rh^{-r}a$ is allowed.
The order of completion is therefore part of the definition.
Equivalently, $\mathcal A$ is the inverse limit of its finite
$t$-jets, each with coefficients in $A((h))$. It is complete
and separated for the $t$-adic filtration. There is no condition
requiring the coefficients of an element of $A((h))$ to
lie in a common finite-dimensional subspace of $A$.

The full Hochschild space of arity $q$ is
$C^q(A,A)=\Hom_\CC(A^{\otimes q},A)$, including $C^0(A,A)=A$.
We use precisely
$C^q(A,A)((h))[[t]]$, with cohomological degree $q-1$.
For a cochain lower bound $B$ and input bounds $b_1,\ldots,b_q$,
the coefficient at $h^d$ is a finite sum over nonnegative
indices with total $d-B-\sum_i b_i$; its output lower bound is
$B+\sum_i b_i$. These estimates justify evaluation and composition
in these spaces, without enlarging them to arbitrary maps on
completed functions.

Finally,
\[
 \log T=\sum_{s\geq1}\frac{(-1)^{s+1}}s(T-\Id)^s
 \in t\mathcal E.
\]
At outer order $r$, only $s\leq r$ occur, with finitely many
positive compositions of $r$; their Laurent bounds add to $-rC$.
For $1\leq s\leq r$ there are exactly
$\binom{r-1}{s-1}$ such compositions, so this is a finite sum of
Laurent series at every outer order. The inverse series obeys
the same lower bound, by the identical calculation without
the rational logarithm coefficients.
Thus the logarithm is a legitimate full Hochschild gauge operator.
It is not claimed to be a derivation of one fixed product.
The next section applies formality to the resulting associative
comparison.
